\documentclass[11pt]{article}
\usepackage{amsmath,amssymb,amsthm,hyperref}
\newtheorem{theorem}{Theorem}
\newtheorem{lemma}{Lemma}
\newcommand{\E}{\mathbb E}
\title{A qualitative strengthening of the three-halves endpoint bound}
\author{Research proof: independently reviewed}
\date{30 September 2026}
\begin{document}\maketitle
\noindent\textbf{Status.} Complete proof, independently reviewed by the root and rational-design
research agents. Its endpoint-profile and variance-sensitive Laguerre
inputs have also been independently reviewed. This qualitative result is
superseded by the stronger near-quadratic bootstrap; it is retained as
a separate structural argument. No explicit exponent exceeding $3/2$
is claimed in this note.

Let a sequence of monotone uniform-interval comparison tensors have
$m\to\infty$ columns, arbitrary positive row weights, and
$\E_w\prod_{j\in S}p_j=1/(|S|+1)$. Put $u_j=1-p_j$, $s=\sum_j u_j$,
$\ell=\lfloor m/4\rfloor$, $N=m-\ell$, $\alpha=\ell/m$, and
$z_j=\ell u_j$. Define the row range
\[
 H_m=\max_j z_j-\min_j z_j.
\]
The previously reviewed uniform endpoint-profile theorem, combined
with finite mean-deficit balance, gives
\begin{equation}\label{profile}
 \sup_{s\le L}H_m\longrightarrow0\qquad(L<\infty\text{ fixed}).
\end{equation}
Indeed, for $s\le L$ the finite bound gives $\max u_j=O_L(1/m)$,
so all entries are positive and $\lambda=\sum u_j/(1-u_j)=s+O_L(1/m)$.
The endpoint profile then states uniformly that $m u_j=s+o(1)$.
The same finite bound gives globally
\begin{equation}\label{global}
 H_m\le C(s+1),\qquad
 \E_w(s+1)^a e^{-cs}\le C_{a,c}/m
\end{equation}
for every fixed $a\ge0,c>0$.

We use the mixture $\mu(q)=(\ell+1)w_q\E_B A_B(q)$,
$A_B=\prod_{j\in B}p_j$, where $B$ is a uniform $\ell$-subset, and
$x=\alpha s$. Exact polarization expectations equal the beta law
$b_\ell(z)=(\ell+1)\ell^{-1}(1-z/\ell)^\ell$ on $[0,\ell]$.

\begin{lemma}[Range-sensitive comparison]\label{compare}
Fix $\kappa>0$ and put $d=\lfloor\kappa\sqrt m\rfloor$.
Let $Q_a(z)=\sum_{j=0}^d c_{aj}L_j(\theta z)$ for a fixed sufficiently
small $\theta>0$, and let $B_a=\sum_j|c_{aj}|$, $B=B_1B_2$.
Uniformly over these coefficients,
\[
 |\E_\mu Q_1(x)Q_2(x)-\beta_\ell(Q_1Q_2)|
 =o(Bd/m).
\tag{1}
\]
\end{lemma}
\begin{proof}
We detail why each formerly $O(Bd/m)$ error acquires a little-oh factor.
For a complementary subset $R$, let $v=N^{-1}\sum_{j\in R}z_j$.
The centered variance satisfies the sharper estimate
\[
 \sigma^2=\sum_{j\in R}(z_j-v)^2\le NH_m v.
\tag{2}
\]
To see this, write $a=\min_{j\in R}z_j$ and $b=\max_{j\in R}z_j$.
The variance is at most $(b-v)(v-a)\le H_m v$ times $N$.

The all-complex Cauchy estimate and the off-diagonal Gaussian Laguerre
bound give respectively
\[
 \frac{|e_k(z-v)|}{\binom Nk}
 \le(\sqrt e\,\sigma\sqrt k/N)^k,
 \qquad
 \frac{|(Q_1Q_2)^{(k)}(v)|}{k!}
 \le B e^{\theta v}(2\sqrt{\theta d/v})^k/k!.
\]
At $v=0$ the centered error vanishes identically. For $k=2$, using (2),
the normalized integrated error is at most
\[
 C B\frac d m\,m\E_w[H_m e^{-cs}]=o(Bd/m).
\]
Here \eqref{profile} handles $s\le L$ and \eqref{global} makes the
$s>L$ contribution uniformly small as $L\to\infty$.
For all $k\ge3$ we use the coarser $H_m\le C(s+1)$ and the factorial
weighted-moment bound. As in the variance-sensitive polarization lemma,
their total is
\[
 O\left(B(d/m)^{3/2}\right)=o(Bd/m).
\]
This proves the improved polarization comparison.

For the mixture comparison, work first on $s\le m/256$, where
$p_i\ge1/2$. Under the success-weighted subset law put $\Delta=v-x$.
The exact bias formula and the identity
$\sum_j(u_j-s/m)^2\le (\max u_j-\min u_j)s$ give
\[
 0\le\E\Delta\le C H_m x/m.
\tag{3}
\]
For the variance, the subset size is fixed. Thus subtracting
$a=\min_jz_j$ from every sampled coefficient leaves its variance
unchanged. The remaining coefficients are nonnegative, so the negative
pairwise inclusion covariances imply
\[
 \operatorname{Var}\Delta
 \le N^{-2}\sum_j(z_j-a)^2
 \le N^{-2}H_m\sum_jz_j
 \le C H_m x/m.
\tag{4}
\]
Consequently $\E\Delta^2\le C H_m x/m+C H_m^2x^2/m^2$.

The off-diagonal derivative bounds are
$|P'(x)|\le C B\sqrt{d/x}e^{\theta x}$ and
$|P''(z)|\le C Bd e^{\theta z}/z$ for $P=Q_1Q_2$.
The integral Taylor remainder obeys
\[
 |P(v)-P(x)-P'(x)(v-x)|
 \le C Bd e^{\theta\max(x,v)}\frac{(v-x)^2}{x}
\tag{5}
\]
for $x>0$, including $v=0$. Indeed, with $v=x+\Delta$,
$(1-t)/(x+t\Delta)\le1/x$ for $0\le t<1$ and $v\ge0$.
At $x=0$ all deficits and all errors vanish.
Combining (3)--(5), averaging with $\overline A\le e^{-\alpha s}$,
and normalizing yields a bound of the form
\[
 C B\frac d m\,m\E_w\left[
   (H_m(1+s)+H_m^2(1+s)/m)e^{-cs}\right]
 =o(Bd/m)
\]
by \eqref{profile}--\eqref{global}. The region $s>m/256$ has
exponentially small error by the previously established crude
uniform-subset comparison, even for $d=O_\kappa(\sqrt m)$.
This proves (1).
\end{proof}

\begin{theorem}
Under the tensor hypotheses above,
\[
 m^{3/2}\max\{w_1,w_K\}\longrightarrow0.
\]
In particular the exceptional-die optimum satisfies
$g(n)/n^{3/2}\to\infty$.
\end{theorem}
\begin{proof}
Use the same scaled Laguerre taper
$T_d(y)=\sum_{j<d}(1-j/d)L_j(y)$ and
$R_d(z)=(2/(\theta d)-z)T_d(\theta z)^2$.
Its beta expectation is bounded below by a fixed positive constant.
The recurrence identity
$yT_d=d^{-1}\sum_{j<d}L_j-L_d$ factors $R_d$ into Laguerre sums
with coefficient norms $O(d)$ and $O(1)$. Lemma~\ref{compare} therefore
gives an error $o(d^2/m)=o_\kappa(1)$ for every fixed $\kappa$.
Hence the last row satisfies $x_K<2/(\theta d)$ for all sufficiently
large $m$.

The normalized Christoffel square at this row has coefficient norm
$O(1)$, value one there, and beta integral $O(1/d)$.
Lemma~\ref{compare} gives
\[
 \mu(K)\le C/d+o(d/m)=\frac{C}{\kappa\sqrt m}
                       +o_\kappa(m^{-1/2}).
\]
Also $s_K=O(1/d)$, so every $A_B(K)\ge1-s_K\ge1/2$ eventually.
Thus
\[
 \limsup_{m\to\infty}m^{3/2}w_K\le C/\kappa.
\]
Since $\kappa>0$ was arbitrary, the limit is zero. Reversal and
complementation give the first-row statement. Equal weights give
the asserted consequence for dice.
\end{proof}
\end{document}
